\addcontentsline{toc}{subsection}{b. Ensayos}
\section*{Ensayo por Saúl Vega Navas.}
\addcontentsline{toc}{subsubsection}{Ensayo por Ensayo por Saúl Vega Navas.}

\subsection*{Introducción.}
El principal objetivo del texto es el de explicar la situación contemporánea de nuestros datos personales como parte de la red de datos que compone al Internet, su valor en el mercado digital y su uso por terceros a costa de nuestra privacidad y del control que tenemos sobre ellos en los medios de comunicación en línea, resaltando el hecho de que, como usuarios, por lo general no tenemos una perspectiva clara de su valor.

Personalmente, la importancia de este tema radica en la valoración de la confidencialidad de la información que nos identifica como individuos, no sólo con respecto a la influencia que ésta refleja en el mercado, sino también con respecto al poder de decisión y confidencialidad con los que decidimos hacer uso de ella y al rol que tomamos como usuarios de Internet frente a terceros: un individuo sin conciencia de la importancia de la preservación de su identidad digital se vuelve simplemente un conjunto de datos sin criterio en el uso de sus datos en Internet por cualquiera.
\subsection*{Desarrollo.}
Actualmente, nos encontramos en la cuarta revolución industrial, caracterizada por el acceso a medios no solo de comunicación e información en cualquier momento con una facilidad sin precedentes, sino también de almacenamiento, gestión y procesamiento en escalas masivas. Es por esto que la economía y, por consecuencia, la sociedad en general, ha tomado un nuevo enfoque en el que la forma de almacenar y obtener datos se ha vuelto un factor fundamental en la influencia que una entidad puede ejercer sobre su entorno, como lo es la forma en que una empresa puede aprovechar para su beneficio la eficacia con la que gestiona la información de sus finanzas, de sus ambientes laborales o, evidentemente, de sus activos.

La información personal se ha convertido en uno de los más importantes activos para las empresas de tecnología: con los usuarios de un servicio en línea generando un rastro de datos digital con cada interacción que tienen con la mayoría de dispositivos electrónicos, el consumo de tiempo que éstos invierten en dichas interacciones se vuelve un recurso aprovechable por las empresas que los recopilan.

No obstante, todo esto es a costa de minimizar la privacidad de los consumidores, y llegando a desviar la atención de las empresas, de la comodidad del usuario hacia la forma de recolectar datos con cada vez mayor precisión y, con ello, adquirir una postura creciente en el mercado digital, recurriendo a medidas extraoficiales como lo fue el caso de Cambridge Analytica en 2018, entidad que recurrió a recopilar los datos de hasta 87 millones de perfiles de Facebook para influir en las recomendaciones de los usuarios y, con ello, en sus decisiones de votaciones electorales.\footnote{Amnistía Internacional (2019). \href{https://www.amnesty.org/es/latest/news/2019/07/the-great-hack-facebook-cambridge-analytica/}{El caso Facebook-Cambridge Analytica}}

Es por esto que, concordando con el artículo, si bien es importante el apoyo a propuestas tecnológicas orientadas al control de los datos de los usuarios para sí mismos, considero que nuestro rol digital como usuarios debe sustentarse en la conciencia que tenemos sobre cómo la información que damos a disposición de otros en Internet es valiosa en el mercado, peligrosa si no se concede con precaución (lo que podría derivar en casos de riesgo a la integridad personal como en el robo de datos bancarios o la suplantación de identidad) pero, sobre todo, valiosa para definirnos como individuos dentro y fuera de los medios de comunicación digital.

Por otra parte, como menciona el artículo, los usuarios de los servicios en línea justificamos la inserción de nuestros datos en la confianza que tenemos al actuar en Internet. Por lo general, esta confianza no se sustenta exclusivamente en el criterio que tenemos sobre nuestra conducta al proporcionar datos de ciertas índoles, sino también sobre la capacidad de administrarlos ante consultas y ante ataques de terceros; además, el almacenamiento de datos no es con el único propósito de comercializar con ellos, sino también con el de tener una nueva perspectiva sobre la calidad de sus servicios.

Es por esto que, dado el valor sensible y las restricciones sobre los datos identificables, es fundamental que las partes involucradas en el diseño y mantenimiento de bases de datos estructuren esquemas que cumplan con la triada de ciberseguridad, y es que los datos deben ser confidenciales, íntegros y accesibles (sólo a quienes estén autorizados), por lo que desarrollar un diseño óptimo para las consultas gestionables y eficaces debe ser una parte esencial del estudio de los fundamentos de las bases de datos.

Asimismo, considero que este planteamiento requiere ser extendido hacia otras áreas de la tecnología para beneficio de los consumidores, manteniendo una simetría en el control informático que existe entre la plataforma y el usuario con respecto a sus datos personales, como lo es principalmente el desarrollo del Internet de las Cosas y los sensores, recopilando información útil no solo para las plataformas, sino también para las personas de quienes la obtienen, como lo es el proyecto de ``Hub of All Things'', mencionado en el texto. 

\subsection*{Conclusión.}

La defensa de nuestra identidad digital y el control irrestricto sobre nuestros datos personales no es únicamente una cuestión de privacidad mercantil, sino de identidad como personas, por lo que es importante consolidar la soberanía de nuestros datos personales con respecto a nuestros estándares de bienestar. Además, con la incesante expansión del Internet de las Cosas y la creciente presencia de sensores en nuestro día a día, es evidente que la recolección de datos es cada vez más exhaustiva, lo que presenta un desafío técnico y ético en la integración de la seguridad y la privacidad desde el diseño, revirtiendo así la asimetría de poder actual para que el usuario siempre tenga la última palabra sobre el destino de su huella digital.

Si bien las implicaciones de cómo esto será abordado son contundentes, el fomentar la adopción de modelos centrados en el usuario permitirá comenzar a desarrollar una economía digital sostenible y en armonía con la dignidad, la ciberseguridad y derechos como personas en la red.
